\section{Method}

\begin{itemize}
    \item{software}
    \item{choice of sim box and amoeba parameters}
    \item{implementation of study conditions}
    \item{parallel array}
    \item{optimisation choices}
    \item{number of timesteps -- lament the relative brevity $\mathcal{O}(10^{-9}s)$ of md sims}
    \item{equilibration -- show plot}
\end{itemize}

There are many different molecular dynamics engines on the market such as GROMACS and some others
whose names I cannot quite remember.
The vehicule of choice for these ice simulations however is LAMMPS (Large-scale Atomic/Molecular Massively
Parallel Simulator).
It's choice was motivated by two factors: it is highly respected, with native support for OpenMPI and 
OpenMP and exhaustive documentation, and it recently implemented the AMOEBA force field (which was 
originally only available on Ponder Lab's own software, Tinker MD).
This is great for me, a complete novice at molecular dynamics!
Now, AMOEBA water has its properties listed in a Tinker '.prm' file with some select properties (which
are of particular importance surely) in Table something along with some of the same parameters
listed for other models of water and some experimental data.
The ice structures were generated using a program called GenIce2 (which I'm not sure if it is relevant 
to mention) and then loaded into LAMMPS.
For Ice--VII there were 2000 molecules and for Ice--VIII there were 1728.
There is certainly something intelligent to say here about how this number is sufficient but also necessary
for observing long-range coordination effects.
There is also room for a couple of figures demonstrating the unit cells potentially if this wasn't already 
in the introduction.
The important thing here is periodic boundary conditions.
We know that a real periodic crystal extends in all three spatial dimensions for many orders of magnitude
worth of atoms, but here we only have 5 or so molecules on either side of the ones in the middle.
We also know that crystal grain boundaries, defects etc. do not behave like the rest of the crystal.
How then can we be sure that the molecules inside our simulation box will behave as if they are in a
real crystal?
We just pretend that they are! 
Periodic boundary conditions can have a massive effect: I'm sure that there is a certain amount of proof
that I could insert right here.
Implementing molecular dynamics simulations can be a bit of an art.
Maintaining the stability of a system with so many different moving parts can be difficult.
The energy required in shortening the covalent bond in water by 10\% is huge and there is probably a number
that I can quote right here.
Hence system setup and conditions are important.
If, say, one molecule were placed on top of another, the Pauli and electrostatic repulsions would send the two
flying off into space immediately.
The same thing can be said with the initial velocities because if they wiggle too much at the start bad things
might happend
Care also needs to be taken with adjusting macroscopic quantities like temperature, pressure and volume.
Adjustments that are too rapid and/or aggresive may destabilize the system, with molecules potenially ending 
up too close to each other or something.
Changing things too slowly eats up computer time.
Furthemrore, since MD siulations preserve the total energy not the temperature or something like that, small
adjustments need to be made constantly to ensure that the simulation conditions are consistent.
Fortunately, people have come up with increasingly clever ways to do this.
We are doing NPT simulations here, and since the number of particles is always constant in a simulation like 
this, we are controlling the pressure and the temperature.
The tools to be implemented are a thermostat and a barostat: spceifically the Nosé-Hoover thermostat and the
Parinello-Rahman barostat.
The math behind these and a short description of what they actually do can be put here.
Potentially could include a plot showing how the temperature and pressure actually oscillate.
Potentially one of the more important things to mention here is the importance of equilibration,
and how the oscillations of the controls can mess with results.
The equilibration period chosen is 250 ps which is pretty standard but also there is some evidence 
from my simulations that coudld justify this choice.
The actual implementation of the procedure is shown in Figure \ref{fig:flowchart}.

%
%         Code diagram
%

\tikzstyle{startstop} = [
    rectangle, 
    rounded corners, 
    minimum width=3cm, 
    minimum height=1cm,
    text centered, 
    draw=black, 
    fill=red!30,
]
\tikzstyle{io} = [
    trapezium, 
    trapezium left angle=70, 
    trapezium right angle=110, 
    minimum width=3cm, 
    minimum height=1cm, 
    text centered, 
    draw=black, 
    fill=blue!30,
]
\tikzstyle{process} = [
    rectangle, 
    minimum width=3cm, 
    minimum height=1cm, 
    text centered, 
    draw=black, 
    fill=orange!30,
]
\tikzstyle{decision} = [
    diamond, 
    minimum width=3cm, 
    minimum height=1cm, 
    text centered, 
    draw=black, 
    fill=green!30,
]
\tikzstyle{arrow} = [
    thick,
    ->,
    >=stealth,
]

\begin{figure}
    \centering
    \begin{tikzpicture}[node distance=2cm]
        \node (start) [startstop] {Start};
        \node (in1) [io, below of=start] {Input};
        \node (pro1) [process, below of=in1] {NVT};
        \node (pro2) [process, below of=pro1] {NPT ramp};
        \node (pro3) [process, below of=pro2] {NPT aniso};
        \node (out1) [io, below of=pro3] {thermo data, dump files, restart};
        \node (stop) [startstop, below of=out1] {Done!};

        \draw [arrow] (start) -- (in1);
        \draw [arrow] (in1) -- (pro1);
        \draw [arrow] (pro1) -- (pro2);
        \draw [arrow] (pro2) -- (pro3);
        \draw [arrow] (pro3) -- (out1);
        \draw [arrow] (out1) -- (stop);
    \end{tikzpicture}
    \caption{Code flowchart}
    \label{fig:flowchart}
\end{figure}


Allowing the simulation to first equilibrate in the NVT ensemble before adjusting the pressure
was much more stable.
Using the anisotropic NPT ensemble for the real thing was motivated mainly by the different known
unit cell structures of the two phases. 
Allowing the crystal to bend and stretch was ideally to reduce suppression of phase transitions. \\
The basic idea was to do simulation runs of the two crystal phases in the same temperature range 
around the phase boundary at a constant pressure.
A `cooling run' for Ice--VII and a `heating run' for Ice--VIII, seeing if one (or both) would change
into the other.
I should drop some numbers here and with reference to the phase diagram. 
To start: these were done in the 3 GPa regime, where the VII/VIII boundary is approximately a straight
line.
The temperature ranged from 230-520K with increments of 10K. 
Although the actual phase boundary lies at around 275K at these conditions, as is typical for molecular
dynamics simulations, a certain amount of hysteresis is expected and the temperatures were required to be
kicked up a notch
Furthermore, considering that the VII $\to$ VIII transition is a disorder $\to$ order transition, the
hysteresis was expected to be more significant and the focus was placed on producing the easier VIII $\to$
VII transition.
Because MD takes forever on the computer, and the equilibration period at each temperature is the
same 250 ps (which was taking about 6hr), doing a temperature ramp \textit{i.e.} heating up the
same system slowly was impractical.
The best thing to do was to run simultaneous simulations at different temperatures and inspect the final
states of each.
These parallel arrays had an average performance of 0.9819 ns day$^{-1}$ (fake number) and were run 
for 4000000 timesteps, achieving a total real time of 4 ns.
Certain optimization choices were made, including choosing an Ewald grid based on numbers which were
multiples of 8 ($2^3$) and compiling the LAMMPS executable with the AMD compiler suite for the 
Slurm node architecture.

